\documentclass[12pt, a4paper]{article}
%--------------------------------------------------------
\input{pacotes} %Não alterar
\input{cabecalho}
%--------------------------------------------------------
\begin{document}
%--------------------------------------------------------
\thispagestyle{firstpage} % Aplica o estilo fancy apenas na primeira página
\vspace*{0.2cm}
%--------------------------------------------------------
\begin{center}
    \singlespacing
    \textbf{TÍTULO (FONTE ARIAL, TAMANHO 12, CAIXA ALTA, NEGRITO, CENTRALIZADO, ESPAÇAMENTO ENTRE LINHAS SIMPLES)}
\end{center}\vspace{-2.em}
%--------------------------------------------------------
\begin{center}
    \singlespacing
    \textit{\textbf{TÍTULO EM INGLÊS (FONTE ARIAL, TAMANHO 12, CAIXA ALTA, NEGRITO, ITÁLICO, CENTRALIZADO, ESPAÇAMENTO ENTRE LINHAS SIMPLES)}}
\end{center}\vspace{-2.em}
%--------------------------------------------------------
\begin{flushright}
    \singlespacing
    \textbf{Nome do autor 1.}\footnote{Vínculo Institucional. E-mail}\\
    \textbf{Nome do autor 2.}\footnote{Vínculo Institucional. E-mail}\\
    \textbf{Nome do autor 3.}\footnote{Vínculo Institucional. E-mail}\\
    \textbf{Nome do autor 4.}\footnote{Vínculo Institucional. E-mail}\\
    \textbf{Nome do autor 5.}\footnote{Vínculo Institucional. E-mail}
\end{flushright}\vspace{-.5em}
%--------------------------------------------------------
{\setlength{\parindent}{0cm} \singlespacing %tamanho do parágrafo
%--------------------------------------------------------
\textbf{RESUMO (FONTE ARIAL, TAMANHO 12, CAIXA ALTA, NEGRITO, JUSTIFICADO, ESPAÇAMENTO ENTRE LINHAS SIMPLES)}

O resumo é elemento obrigatório e deve ser elaborado conforme a ABNT-NBR 6028:200. Deve ser escrito em dois idiomas: português e inglês. A sequência dos resumos está condicionada ao idioma no qual o artigo está escrito: artigo redigido em português (manter o padrão do \textit{template}); artigo redigido em inglês (resumo em inglês > resumo em português). O texto do resumo deve usar fonte Arial, tamanho 12, alinhamento de parágrafo justificado, sem recuos à direita ou à esquerda e espaçamento entre linhas simples. Conforme ABNT NBR 6028, o resumo deve conter \textbf{entre 100 e 250 palavras}, ser do \textbf{tipo informativo} e corresponder a uma \textbf{síntese precisa do artigo}, contendo \textbf{obrigatoriamente} introdução, objetivo, matriz teórica, metodologia, resultados e conclusão.\\

\textbf{PALAVRAS-CHAVE:} Use entre \textbf{3 (três) e 5 (cinco)} palavras-chave, separando-as por \textbf{ponto e vírgula} e finalizadas por \textbf{ponto}. Letras iniciais minúsculas, exceto substantivos próprios e nomes científicos.\\
%--------------------------------------------------------

\textbf{ABSTRACT (FONTE ARIAL, TAMANHO 12, CAIXA ALTA, NEGRITO, JUSTIFICADO, ESPAÇAMENTO ENTRE LINHAS SIMPLES)}

Realizar tradução profissional, utilizando linguagem da norma culta e acadêmica. Atenção quanto ao uso de programas de tradução.\\

\textbf{KEYWORDS:} xxxxxxxxxxxxxxxxxxxxxxxxxxxxxxxxxxxxxx

}
%--------------------------------------------------------
\pagestyle{otherpages}
%-------------------------------------------------------
\section{INTRODUÇÃO}
Este é o padrão para a submissão de trabalhos da Categoria 2 da 21ª SNCT, destinados à divulgação de \textbf{pesquisas com resultados parciais ou conclusivos}. Nesta categoria, os trabalhos devem ser submetidos em Português ou Inglês, em forma de resumo expandido com no mínimo 2 e no máximo 4 páginas, incluindo-se as referências bibliográficas. Os trabalhos submetidos que não estiverem de acordo com o formato apresentado por esse padrão \textcolor{red}{serão rejeitados} pelo Comitê Editorial do evento, sem análise do mérito científico.

A \textbf{INTRODUÇÃO} deve contextualizar o objeto de estudo, contendo questões ou problemas de pesquisa, hipóteses (quando houver) e objetivos. Não deve conter ilustrações e tabelas. Não deve ser dividida em seções de 2º e 3º nível.

As citações no texto devem seguir a \textbf{ABNT NBR 10520}. \textbf{Atenção quanto à atualização da norma de citação ocorrida no ano de 2023}.

Saltar uma linha antes e depois de citação direta com mais de 3 linhas, utilizando fonte Arial, tamanho 10, alinhamento justificado e espaçamento entre linhas simples.

Cada citação no corpo do texto deve corresponder a uma referência completa na lista de referências ao final do trabalho. As citações devem ser feitas na língua do artigo. No caso de citações cuja fonte original seja em língua estrangeira, a pessoa autora deve traduzir e indicar na referência a expressão \textbf{“tradução nossa”}.

\begin{enumerate}[noitemsep]
\item[a.] \textbf{Citação indireta} no corpo do texto, com uma pessoa autora: (Sobrenome, ano);
\item[b.] \textbf{Citação indireta} no corpo do texto, com pessoas autoras de obras distintas: (Sobrenome, ano; Sobrenome, ano);
\item[c.] \textbf{ Citação indireta} no corpo do texto, com duas pessoas autoras de uma mesma obra: (Sobrenome; Sobrenome, ano);
\item[d.] \textbf{Citação direta} até três linhas: “Texto da citação direta entre aspas no decorrer do artigo.” (Sobrenome, ano, p. 0);
\item[e.] \textbf{Citação direta} até três linhas com destaque (negrito e/ou itálico) da pessoa autora do artigo: “Texto da \textbf{citação direta} entre aspas no decorrer do artigo.” (Sobrenome, ano, p. 0, grifo nosso);
\item[f.] Usar colchetes para omitir parte de citação direta: “[...] citação direta entre aspas [...] do artigo” (Sobrenome, ano, p. 0);
\item[g.] Citação direta longa (mais de três linhas):
    \begin{adjustwidth}{4cm}{0cm}
        \singlespacing \small
        Utilize fonte Arial, tamanho 10, entrelinha simples, justificado, recuo de parágrafo de 4 cm à esquerda. (Sobrenome, ano, p. 0). Texto texto texto texto texto texto texto texto texto texto texto texto texto texto texto texto texto texto texto texto texto texto texto texto texto texto texto texto texto texto;
    \end{adjustwidth}
       
\item[h.] \textbf{Citação de citação:} usar a expressão latina apud. Na listagem final de referências, inserir a referência da obra acessada, isto é, a obra que citou a referência utilizada, não a obra original. Exemplo: Hurd (1996 apud Leite, 2006) afirma que [...].     
\end{enumerate}

\section{METODOLOGIA (A SEÇÃO PODE SER RENOMEADA E SUBDIVIDIDA DE ACORDO COM A NATUREZA DO TRABALHO)} 

Caracterizar o recorte espacial e temporal da pesquisa quando pertinente. Elencar os fundamentos, caracterizar a pesquisa e apresentar os procedimentos metodológicos de acordo com a natureza do artigo.

\section{RESULTADOS E DISCUSSÃO (A SEÇÃO PODE SER RENOMEADA E SUBDIVIDIDA DE ACORDO COM A NATUREZA DO TRABALHO)}

Apresentar, analisar e discutir os resultados. Comparar com estudos científicos semelhantes quando pertinente. Relacionar os resultados com os fundamentos teóricos.

As ilustrações (desenhos, esquemas, fluxogramas, fotografias, gráficos, mapas, organogramas, plantas, quadros, retratos etc.) e as tabelas devem constar no corpo do texto, não excedendo suas margens. As ilustrações e tabelas devem ser inseridas o mais próximo possível de suas menções no texto. Escrever a denominação da Ilustração ou Tabela com a 1ª letra maiúscula e não abreviar. \textbf{Exemplos:} Na Figura 1 é possível analisar [...]. Conforme observado na Tabela 1 [...].

O título da ilustração ou tabela deve ser inserido na parte superior, fonte Arial, tamanho 12, alinhamento justificado, espaçamento entre linhas simples, numerados sequencialmente. 

\textbf{Exemplos:}

\begin{figure}[H]
    \centering
    \caption{Título da Ilustração.}
    \includegraphics[width=7cm]{SNCT2024.PNG}\\
    \caption*{\footnotesize Fonte: SNCT, 2024.}
        \label{fig:enter-label}
\end{figure}

\begin{table}[H]
    \centering
    \caption{Título da Tabela.}
    \begin{tabular}{ccc}
    \hline
    \multicolumn{1}{c}{\textbf{Negrito}} & \textbf{Negrito} & \textbf{Negrito} \\ \hline
    Texto             & Texto      & Texto            \\
    Texto                                  & Texto            & Texto            \\
    Texto                                  & Texto            & Texto      \vspace{0.2cm}    \\
    \hline
    \end{tabular}
    \caption*{\footnotesize Fonte: SNCT, 2024.}
\end{table} \vspace{-0.8cm} 

Saltar uma linha antes e depois da ilustração ou tabela, utilizando fonte Arial, tamanho 12, alinhamento justificado e espaçamento entre linhas de 1,5 cm.

\section{CONSIDERAÇÕES FINAIS (OU CONCLUSÃO)}
Sintetizar os resultados mais relevantes. Realizar reflexões teóricas que contribuam com o avanço da compreensão do objeto de estudo e campo de investigação. Evidenciar limitações e desafios da pesquisa e do entendimento do objeto de estudo quando pertinente. \textbf{A CONCLUSÃO não deve ser dividida em seções de 2º e 3º nível.}

\section*{REFERÊNCIAS} % Apagar e descolmentar a linha \printbibliography
\vspace{-0.5cm} 
\singlespacing
\begin{flushleft}
\noindent Todas as referências devem ser organizadas conforme a \textbf{ABNT NBR 6023}. Orientações gerais: fonte Arial; tamanho 12; sem recuo de primeira linha; espaçamento entre linhas simples; alinhadas à margem esquerda; separadas entre si por um espaço em branco com espaçamento entre linhas simples; ordem alfabética; hiperlinks desativados; abreviação ou nome completo dos autores (manter o padrão em todas as referências). 
\end{flushleft}

\noindent Exemplos:

\noindent \textcolor{red}{Artigos de periódico ou revista}
\begin{flushleft}
\singlespacing
\noindent 
SOBRENOME, Nome abreviado. \textbf{Título do artigo}. Título da Revista, Local de publicação, número do volume, páginas inicial-final, mês e ano.
\end{flushleft}
\textcolor{red}{Artigos com mais de 3 (três) autores}
\begin{flushleft}
\singlespacing
\noindent 
YAMAKAWA, E. K. et al. Comparativo dos softwares de gerenciamento de referências bibliográficas: Mendeley, EndNote e Zotero. \textbf{Transinformação}, v. 26, p. 167-176, 2014.
\end{flushleft}
\textcolor{red}{Artigos em um evento}
\begin{flushleft}
\singlespacing
\noindent 
SOBRENOME, Nome abreviado. Título do trabalho apresentado. In: TÍTULO DO EVENTO, nº do evento, ano de realização, local (cidade de realização). \textbf{Título do documento} (anais, resumos, etc). Local: Editora, ano de publicação. Páginas inicial-final.
\end{flushleft}
\textcolor{red}{Monografia, dissertação ou tese}
\begin{flushleft}
\singlespacing
\noindent SOBRENOME, Nome abreviado. \textbf{Título}: subtítulo (se houver). Ano de apresentação. Número de folhas ou volumes. Categoria (área de concentração) – Instituição, Local, ano da defesa.
\end{flushleft}
%------------------ Referências ---------------------
% Descometar a linha abaixo para inseri a seção das referências. As instroções de citação estão no arquivo referencias.bib

%\printbibliography[title = REFERÊNCIAS]
%-------------------------------------------------------

\end{document}